% from ega2-1-fr.tex
\oldpage[II]{174}
en vertu de (\hyperref[II.8.5.1.3-fr]{8.5.1.3}) et
(\hyperref[II.8.5.3.2-fr]{8.5.3.2})~; en outre, la restriction de $r$ au
complémentaire $\widehat{C}_X-i_X(\varepsilon_X(X))$ de la section nulle est
un \emph{isomorphisme} sur le complémentaire
$\widehat{C}-i(\varepsilon(Y))$ de la section nulle en vertu de
(\hyperref[II.8.6.2-fr]{8.6.2}). Si on suppose pour simplifier $Y$ affine,
$\mathcal{S}$ de type fini et engendré par $\mathcal{S}_1$, $X$ est projectif
sur $Y$ et $\widehat{C}_X$ projectif sur $X$
(\hyperref[II.5.5.1-fr]{5.5.1}), donc $\widehat{C}_X$ est projectif sur $Y$
(\hyperref[II.5.5.5-fr]{5.5.5, (ii)}), et \emph{a fortiori} sur
$\widehat{C}$ (\hyperref[II.5.5.5-fr]{5.5.5, (v)}). On a ainsi un
$Y$-morphisme projectif $r:\widehat{C}_X\to\widehat{C}$ (dont la restriction
à $C_X$ est un $Y$-morphisme projectif $C_X\to C$) qui \emph{contracte $X$
dans $Y$} et induit un \emph{isomorphisme} quand on le restreint aux
\emph{complémentaires de $X$ et de $Y$}. On a donc une relation entre $C_X$
et $C$, analogue à celle qui a lieu entre un préschéma éclaté et son préschéma
de départ (\hyperref[II.8.1.3-fr]{8.1.3}). Nous allons effectivement montrer
qu'on peut identifier $C_X$ au spectre homogène d'une
$\mathcal{O}_C$-Algèbre graduée.
\end{env}

\begin{env}[8.7.2]
\label{II.8.7.2-fr}
Gardant les notations de (\hyperref[II.8.6.1-fr]{8.6.1}), considérons pour
chaque $n\geqslant 0$, l'Idéal quasi-cohérent
\begin{equation}
\label{II.8.7.2.1-fr}
  \mathcal{S}_{[n]}=\bigoplus_{m\geqslant n}\mathcal{S}_m
\tag{8.7.2.1}
\end{equation}
de la $\mathcal{O}_Y$-Algèbre graduée $\mathcal{S}$. Il est clair que l'on a
\begin{equation}
\label{II.8.7.2.2-fr}
  \mathcal{S}_{[0]}=\mathcal{S},\qquad
  \mathcal{S}_{[n]}\subset\mathcal{S}_{[m]}
  \qquad\text{pour }m\leqslant n
\tag{8.7.2.2}
\end{equation}
\begin{equation}
\label{II.8.7.2.3-fr}
  \mathcal{S}_{[n]}\mathcal{S}_{[m]}\subset\mathcal{S}_{[m+n]}.
\tag{8.7.2.3}
\end{equation}

Considérons le $\mathcal{O}_C$-Module associé à $\mathcal{S}_{[n]}$, qui est un
Idéal quasi-cohérent de $\mathcal{O}_C=\widetilde{\mathcal{S}}$
(\hyperref[II.1.4.4-fr]{1.4.4})
\begin{equation}
\label{II.8.7.2.4-fr}
  \mathcal{I}_n=(\mathcal{S}_{[n]})^{\sim}.
\tag{8.7.2.4}
\end{equation}

On déduit alors de (\hyperref[II.8.7.2.2-fr]{8.7.2.2}) et
(\hyperref[II.8.7.2.3-fr]{8.7.2.3}), utilisant
(\hyperref[II.1.4.4-fr]{1.4.4}) et (\hyperref[II.1.4.8.1-fr]{1.4.8.1}),
les formules analogues
\begin{equation}
\label{II.8.7.2.5-fr}
  \mathcal{I}_0=\mathcal{O}_C,\qquad
  \mathcal{I}_n\subset\mathcal{I}_m
  \qquad\text{pour }m\leqslant n
\tag{8.7.2.5}
\end{equation}
\begin{equation}
\label{II.8.7.2.6-fr}
  \mathcal{I}_n\mathcal{I}_m\subset\mathcal{I}_{n+m}.
\tag{8.7.2.6}
\end{equation}

On est donc dans les conditions de (\hyperref[II.8.1.1-fr]{8.1.1}), ce qui
conduit à introduire la $\mathcal{O}_C$-Algèbre graduée quasi-cohérente
\begin{equation}
\label{II.8.7.2.7-fr}
  \mathcal{S}^{\natural}
  =\bigoplus_{n\geqslant 0}\mathcal{I}_n
  =\left(\bigoplus_{n\geqslant 0}\mathcal{S}_{[n]}\right)^{\sim}.
\tag{8.7.2.7}
\end{equation}
\end{env}

\begin{proposition}[8.7.3]
\label{II.8.7.3-fr}
Il existe un $C$-isomorphisme canonique
\begin{equation}
\label{II.8.7.3.1-fr}
  h:C_X\xrightarrow{\sim}\operatorname{Proj}(\mathcal{S}^{\natural}).
\tag{8.7.3.1}
\end{equation}
\end{proposition}

\begin{proof}
Supposons d'abord $Y=\operatorname{Spec}(A)$ affine, d'où
$\mathcal{S}=\widetilde{S}$, où $S$ est une $A$-algèbre graduée à degrés
positifs et $C=\operatorname{Spec}(S)$. La définition
(\hyperref[II.8.2.6.2-fr]{8.2.7.4}) montre que l'on a alors, avec les
notations de (\hyperref[II.8.2.6-fr]{8.2.6}),
$\mathcal{S}^{\natural}=(S^{\natural})^{\sim}$. Pour définir
(\hyperref[II.8.7.3.1-fr]{8.7.3.1}), considérons un élément homogène
$f\in S_d$ ($d>0$) et l'élément correspondant $f^{\natural}\in S^{\natural}$
(\hyperref[II.8.2.6-fr]{8.2.6})~; le $S$-isomorphisme
(\hyperref[II.8.2.7.3-fr]{8.2.7.3}) définit donc un $C$-isomorphisme
\begin{equation}
\label{II.8.7.3.2-fr}
  \operatorname{Spec}(S_f^{\geqslant})
  \xrightarrow{\sim}
  \operatorname{Spec}(S_{(f^{\natural})}^{\natural}).
\tag{8.7.3.2}
\end{equation}

